\documentclass[cs4size,a4paper,nofonts]{ctexart}
\usepackage[utf8]{inputenc}
\def\tjf{{\tt{田劲锋}}}
\def\titlec{Linux系统基本操作命令}

\usepackage[a4paper,margin=2.2cm]{geometry} % 页面设置
\usepackage[unicode,breaklinks=true,
colorlinks=true,linkcolor=black,anchorcolor=black,citecolor=black,urlcolor=black,
pdftitle={\titlec},pdfauthor={\tjf}]{hyperref}
\usepackage{latexsym,amsmath,amssymb,bm}
\usepackage{graphicx}
\usepackage{xcolor}
\usepackage{fancyvrb}
\DefineShortVerb{\|}
\fvset{frame=single}

\setmainfont{Times New Roman}
\setCJKmainfont[BoldFont={SimHei}]{SimSun}  % 主要字体：宋体、黑体
\setCJKsansfont[BoldFont={STZhongsong}]{STFangsong} % 次要字体：仿宋、中宋
\setCJKmonofont{KFKai} % 等宽字体：楷体

\usepackage{listings} % 排版代码
\lstset{language=C,
  numbers=left,
  numberstyle=\small,
  basicstyle=\tt\color{black},
  commentstyle=\color{gray},
  keywordstyle=\bfseries\color{blue},
  stringstyle=\color{teal},
  showstringspaces=false,
  frame=tb,
  extendedchars=false,
%     % identifierstyle=\sf,
%     breaklines=true,
  texcl=true,
  mathescape=true,
  escapeinside=` `,
%     % texcl=true,  % 为true时，注释内容被当做LaTeX代码来处理
}


\CJKsetecglue{\hspace{0.1em}}
\renewcommand\CJKglue{\hskip -0.3pt plus 0.08\baselineskip}
\frenchspacing
\widowpenalty=10000
\linespread{1.2} % 行距

\CTEXsetup[name={实验,},number={\arabic{part}}]{part}
\CTEXsetup[name={,.},format={\large}]{section}

\usepackage[inline]{enumitem} % 调整列表样式
\setlist{noitemsep}
\setlist[itemize]{topsep=0pt,partopsep=0pt,itemsep=0pt,parsep=0pt}
\setlist[enumerate]{topsep=0pt,partopsep=0pt,itemsep=0pt,parsep=0pt}
\setlist[enumerate,1]{label={\arabic*.}}
\setlist[enumerate,2]{label={(\arabic*)}}

%\makeindex
\pagestyle{plain}

\newcommand{\image}[3][height=10cm]{%\begin{minipage}[t]{0.5\textwidth}
    \centering
    \includegraphics[#1]{#2}
    \caption{#3}
    \label{fig:#3}
%\end{minipage}
}


\begin{document}

%%%% 开始 %%%%

\setcounter{part}{0}
\def\today{2015~年~5~月~15~日}
\begin{titlepage}

\begin{center}

\includegraphics[height=1cm]{image/haut.png}

\vspace*{1cm}
{\liti\fontsize{48pt}{50pt}{课\quad 程\quad 设\quad 计}}

\vspace*{4cm}
{\fontsize{36}{80}\sf\bfseries \titlec}

\vspace*{1cm}
{\fontsize{30}{70}\sf\bfseries \titlee}

\vfill
{\large
\newcommand{\ctline}[2]{\makebox[6em][s]{\bf #1}：\underline{\makebox[14em][c]{\qquad #2\qquad}}\\}
\ctline{课程设计名称}{数据结构课程设计}
\ctline{专业班级}{计算机 1303 班}
\ctline{学生姓名}{\tjf}
\ctline{学号}{201316920311}
\ctline{指导教师}{白\quad 浩}
\ctline{课程设计时间}{\today}
}

\end{center}

\end{titlepage}


\iffalse
\section{实验目的}
\begin{enumerate}
\item 熟悉Linux操作系统的基本操作命令；
\item 通过在Linux环境下对进程的基本操作，认识进程并区分与程序的区别。
\end{enumerate}

\section{实验环境}
一台装有Linux操作系统（Fedora 7），至少具有256M内存的微机。

\section{预备知识}
\begin{enumerate}
\item vi编辑器的使用
\item 系统管理命令\\
login, logout, man
\item 文件操作命令\\
ls, cat, cp, mv , rm,
\item 目录操作命令\\
cd, mkdir, rmdir, pwd
\item 进程操作命令\\
ps, kill, pstree
\end{enumerate}
\fi

\section{实验内容}
\begin{enumerate}[label={(\arabic*)}]
\item 用vi编辑器建立一个文件，用文件操作命令对该文件进行操作；
\item 使用目录操作命令进行创建目录、改变当前目录、删除目录、查看当前所在目录等操作；
\item 使用进程操作命令查看系统中的进程状态，杀死进程等；
\item 使用man命令了解各个命令的使用方法及参数。
\end{enumerate}

\section{实验要求}
\begin{enumerate}
\item 将各个命令行使用规格写入实验报告；
\item 将各命令操作过程及结果写入实验报告。
\end{enumerate}

\section{实验步骤}

\begin{enumerate}[label={(\arabic*)}]

\item 这里我使用Vim而不是|vi|，在终端中键入|vim|以启动：

\begin{Verbatim}
$ vim ~/.vimrc
\end{Verbatim}

\begin{figure}[htp]
\image{exp1/1-vi.png}{Vim编辑器}
\end{figure}

如图\ref{fig:Vim编辑器}，这里编辑的是Vim的配置文件。

Vim常用三个模式：命令模式、插入模式、可视模式。

在命令模式下，用|h j k l|左下上右移动光标，|w e b|移动到单词首尾，\verb+0 ^ $+移动到行首尾，|gg G|移动到文档首尾。

|yy|复制行，|yw|复制单词，|p P|粘贴，|dd dw D x|剪切。命令前加数字可以重复。

进入插入模式用|i|，或者|a|从光标后插入，|o O|另起新行，|I A|插入到行首尾。

用|Esc|退出插入模式。

可视模式不常用，按|v|进入，可以选择文本。

冒号指令|:w|保存，|:e|编辑，|:q|退出，|:x|保存并退出。

|/pattern|查找，\verb+:%s/find/replace/g+替换。

随时可以参考Vim自带的文档（有中文版），剩下的就是熟能生巧了。

\begin{figure}[htp]
\image{exp1/2-dir.png}{目录操作}
\end{figure}

\clearpage

\item 如图\ref{fig:目录操作}，使用如下命令建立、进入、删除了一个名为test的目录：

\begin{Verbatim}
$ ls
$ mkdir test
$ cd test
$ cd ..
$ rmdir test
$ ls
\end{Verbatim}

这几个命令都有着共同的特点，就是可以跟一个目录名参数。

\begin{itemize}
\item |ls|列出目录内容，常用|-a|列出全部文件，用|-l|列出列表，用|-1|每行单列通过管道传送到|xargs|，新版本的有|-h|可以以人类可读方式来显示文件大小。
\item |mkdir|创建目录，常用|-p|递归创建目录层级。
\item |cd|改变当前目录，是Shell的内建命令。常用|.|为当前目录，|..|为上级目录。
\item |rmdir|移除目录，只能移除空目录，常用|-p|递归移除目录。
\end{itemize}

\begin{figure}[htp]
\image{exp1/3-ps.png}{进程操作}
\end{figure}

\item 如图\ref{fig:进程操作}，键入
\begin{Verbatim}
$ cat &
[1] 1148
[1]  + 1148 suspended (tty input)  cat
\end{Verbatim}
执行一个由标准输入阻塞的|cat|进程，用
\begin{Verbatim}
$ ps
  PID TTY           TIME CMD
  820 ttys000    0:00.61 -zsh
 1148 ttys000    0:00.00 cat
\end{Verbatim}
可以看到在当前终端会话中运行的进程，这里看到进程号为1148，那么执行
\begin{Verbatim}
$ kill 1148
[1]  + 1148 terminated  cat
\end{Verbatim}
就可以结束掉该进程。

\item 经典的UNIX的手册共有七部分，分别是
\begin{enumerate}[label={\arabic*.}]
\item 命令
\item 系统调用
\item 库函数
\item 特殊文件和设备
\item 文件格式
\item 游戏和杂项
\item 启动和登录进程
\end{enumerate}
后来的版本则会增加新的部分，但都大同小异。

使用|man|来查阅手册，可以直接跟命令名，也可以指定其所在部分。如图\ref{fig:man命令}是执行
\begin{Verbatim}
$ man ps
\end{Verbatim}
出现的手册页，用|d p|可以上下翻页，|q|退出。

\begin{figure}[htp]
\image{exp1/4-man.png}{man命令}
\end{figure}

也可以指定其所在章节，如下两行效果是不同的
\begin{Verbatim}
$ man printf
$ man 3 printf
\end{Verbatim}
第一个（如图\ref{fig:man printf}）会显示第一部分的内建指令|printf|，而第二个（如图\ref{fig:man 3 printf}）会显示第三部分C语言|<stdio.h>|提供的|printf()|函数。

\begin{figure}[htp]
\begin{minipage}[t]{0.5\textwidth}\image[width=\textwidth]{exp1/man1.png}{man printf}\end{minipage}
\begin{minipage}[t]{0.5\textwidth}\image[width=\textwidth]{exp1/man2.png}{man 3 printf}\end{minipage}
\end{figure}

GNU计划提供的|info|也是非常好用的手册工具。

应用程序提供的|--help|选项通常也是有效的参考。

\end{enumerate}

\end{document}
